\usepackage{fontspec}
\setmainfont{LibertinusSerif-Regular.otf}[
  BoldFont=LibertinusSerif-Bold.otf,
  ItalicFont=LibertinusSerif-Italic.otf,
  BoldItalicFont=LibertinusSerif-BoldItalic.otf]
\setsansfont{LibertinusSans-Regular.otf}[
  BoldFont=LibertinusSans-Bold.otf,
  ItalicFont=LibertinusSans-Italic.otf]
\setmonofont{DejaVu Sans Mono}[Scale=MatchLowercase]
\usepackage{unicode-math}
\setmathfont{LibertinusMath-Regular.otf}
\usepackage[margin=25mm,headheight=16pt]{geometry}
\usepackage{microtype}
\usepackage{array,booktabs,longtable,calc}
\usepackage{footnotehyper}
\usepackage{xcolor,listings}
\lstset{basicstyle=\small\ttfamily,columns=fullflexible,
  keepspaces=true,showstringspaces=false,breaklines=true,
  breakatwhitespace=false,aboveskip=0.7\baselineskip,
  belowskip=0.7\baselineskip}
% Context titles are typeset separately, never copied as terminal input.
\lstnewenvironment{commandexample}[2][]
  {\lstset{#1,title={\sffamily\bfseries #2}}}{}
% Short operational cautions stay together, without extra package dependencies.
\newenvironment{operationnote}[1]
  {\par\medskip\setlength{\parfillskip}{0pt plus 1fil}%
   \noindent\begin{minipage}{\linewidth}
   \hrule\smallskip{\sffamily\bfseries #1}\par\smallskip\ignorespaces}
  {\par\smallskip\hrule\end{minipage}\par\medskip}
% Keep short procedures and their cautions on one page.
\newenvironment{procedureblock}
  {\par\setlength{\parfillskip}{0pt plus 1fil}%
   \noindent\begin{minipage}{\linewidth}}
  {\end{minipage}\par}
\providecommand{\passthrough}[1]{#1}
\providecommand{\tightlist}{\setlength{\itemsep}{0pt}\setlength{\parskip}{0pt}}
\usepackage{xurl}
\usepackage[unicode,colorlinks=true,linkcolor=blue!45!black,
  urlcolor=blue!45!black,pdftitle={SIM4000 User Manual},
  pdfauthor={SIM4000 project}]{hyperref}
\usepackage{bookmark}
\usepackage{makeidx}
\makeindex
% Unlike a web link, this reference is useful in print as well as in the PDF.
\newcommand{\manualref}[2]{\hyperref[#1]{#2 (p.~\pageref*{#1})}}
\usepackage[automark]{scrlayer-scrpage}
\clearpairofpagestyles
\ihead{SIM4000 User Manual}
\ohead{\ifnum\value{chapter}>0 Chapter~\thechapter\fi}
\ofoot*{\pagemark}
\setcounter{tocdepth}{1}
\setcounter{secnumdepth}{1}
\setlength{\emergencystretch}{3em}
